%% This document is created by 
%%  Dr. Putu Harry Gunawan
%% and edited by 
%%  Gia S. Wulandari (2023)
%% Template untuk Proposal TA 1 dan TA
%% Template ini digunakan untuk penulisan proposal TA 1 atau TA Fakultas Informatika, Telkom University.
%%  New Update 28 Pebruari 2017 : Bookmarks pdf, Pembimbing hanya satu

\documentclass[a4paper,12pt,oneside]{book}
\usepackage[utf8]{inputenc}
\usepackage{fix-cm} % Allow Computer Modern at arbitrary font sizes
\usepackage{sectsty}
\usepackage{graphicx}
\usepackage{epstopdf}
\usepackage{algorithm}
\usepackage{algpseudocode}
\usepackage{array}
\usepackage[table]{xcolor}
\usepackage{anysize}
\usepackage{amsmath}
\usepackage{amssymb}
\usepackage[bahasa]{babel}
\usepackage{indentfirst} %Spasi untuk paragraf pertama
\usepackage{geometry}
\usepackage{multirow}% http://ctan.org/pkg/multirow
\usepackage{hhline}% http://ctan.org/pkg/hhline
\marginsize{4cm}{3cm}{3cm}{3cm} %{left}{right}{top}{bottom}
\usepackage[compact]{titlesec} 
\usepackage{etoolbox}
\usepackage[numbers]{natbib} %for Harvard syle
\usepackage[hidelinks,bookmarks,hypertexnames=false]{hyperref}
%\usepackage[pageanchor]{hyperref}
\usepackage{bookmark}


\makeatletter
\patchcmd{\ttlh@hang}{\parindent\z@}{\parindent\z@\leavevmode}{}{}
\patchcmd{\ttlh@hang}{\noindent}{}{}{}
\makeatother

\chapterfont{\centering}
\newcommand{\bigsize}{\fontsize{16pt}{14pt}\selectfont}
\chapterfont{\centering\bigsize\bfseries}
\sectionfont{\large\bfseries}
\usepackage{tikz}
\usetikzlibrary{shapes.geometric, arrows}
%\renewcommand{\chaptertitle}{BAB}
\renewcommand{\thechapter}{\Roman{chapter}}
\renewcommand\thesection{\arabic{chapter}.\arabic{section}}
\renewcommand\thesubsection{\thesection.\arabic{subsection}}
\renewcommand{\theequation}{\arabic{chapter}.\arabic{equation}}
\renewcommand{\thefigure}{\arabic{chapter}.\arabic{figure}}
\renewcommand{\thetable}{\arabic{chapter}.\arabic{table}}

\renewcommand\bibname{Daftar Pustaka}
\addto{\captionsbahasa}{\renewcommand{\bibname}{Daftar Pustaka}}
\usepackage{fancyhdr}
\pagestyle{fancy}
\lhead{}
\chead{}
\rhead{}
\lfoot{}
\cfoot{\thepage}
\rfoot{}
\renewcommand{\headrulewidth}{0pt}

\makeatletter

%%%%%%%%%%%%%%%%%%%%%%%%%%%%%%%%%%%%%%%%%%%%%%%%%%%%%%%%%%%%
%
%  Berikut adalah data-data yang wajib diisi oleh mahasiswa
%
%%%%%%%%%%%%%%%%%%%%%%%%%%%%%%%%%%%%%%%%%%%%%%%%%%%%%%%%%%%%

\title{Analisis Trade-off antara Akurasi dan Efisiensi Komputasi dalam Deteksi Sarkasme Menggunakan DistilBERT dan RoBERTa}\let\Title\@title   %Judul dalam bahasa Indonesia

\newcommand{\EngTitle}{Trade-off Analysis Between Accuracy and Computational Efficiency in Sarcasm Detection Using DistilBERT and RoBERTa}  %Judul dalam bahasa Inggris

\author{Falah Razan Hibrizi}  \let\Author\@author  %Nama mhs
\newcommand{\NIM}{103012300245} %Jangan lupa hapus tanda ( dan )
\newcommand{\Prodi}{Informatika}
\newcommand{\KK}{(Data Science and Intelligent System)} %UNTUK TA, Jangan lupa hapus tanda ( dan )
\newcommand{\Gelar}{S.Kom.} % UNTUK TA
\date{\the\year}           \let\Date\@date %Masukkan hanya tahun saja
\newcommand{\Tanggal}{\the\day} % Tanggal Pengesahan
\newcommand{\Bulan}{Oktober} % Bulan Pengesahan
\newcommand{\PembimbingSatu}{(Dr. Calon Pembimbing 1, M.Kom)}
\newcommand{\NIPPembimbingSatu}{123456}
\newcommand{\PembimbingDua}{(Dr. Calon Pembimbing 2, M.Kom)}
\newcommand{\NIPPembimbingDua}{123456}
\newcommand{\Kaprodi}{Dr. Erwin Budi Setiawan, S.Si., M.T.}
\newcommand{\NIPKaprodi}{00760045}
\newif\ifPembimbingHanyaSatu

%%%% WARNING kode berikut ini diaktifkan Jika pembimbing hanya satu %%%%

%\PembimbingHanyaSatutrue

%%%%%%%%%%%%%%%%%%%%%%%%%%%%%%%%%%%%%%%%%%%%%%%%%%%%%%%%%%%%%%%%%%

\makeatother
\linespread{1}


\begin{document}
%%%%%%%%%%%%%%%%%%%%%%%%%%%%%%%%%%%%%%%%%%%%%%%%%%%%%%%%%%
% Dimulai dengan cover, lembar persetujuan, Abstrak, dll
%%%%%%%%%%%%%%%%%%%%%%%%%%%%%%%%%%%%%%%%%%%%%%%%%%%%%%%%%%%
\pagenumbering{roman} 
\begin{titlepage}
\thispagestyle{empty}
%\vspace*{0.7cm}
\pdfbookmark{Cover}{ }
{\centering
\large
{\bigsize\bf \Title}\\
\vspace{ 3cm}
\rm
\textbf{Proposal Tugas Akhir}\\
\vspace{0.5 cm}
\textbf{IF-47-DSIS.05}\\
\vspace{0.5 cm}
 \textbf{\NIM}\\ \textbf{\Author}\\

\vspace{1.5 cm}

\begin{figure}[h]
{\centering {\includegraphics[scale=0.17]{Tel-U-Logo}}\par}
\end{figure}

\vspace{2 cm}
{\bigsize\textbf{Program Studi Sarjana \Prodi}\\
\vspace{0.5 cm}
\textbf{Fakultas Informatika}\\
\vspace{0.5 cm}
\textbf{Universitas Telkom}\\
\vspace{0.5 cm}
\textbf{Bandung}\\
\vspace{0.5 cm}
\textbf{\Date}\\}
}
\pagebreak
\thispagestyle{empty}
\pdfbookmark{Lembar Persetujuan}{ }
{\centering
\textbf{\large Lembar Persetujuan}\\
\vspace{0.5cm}

\textbf{\Title}\\
\vspace{0.5cm}

\textbf{\textit{\EngTitle}}\\
\vspace{0.5cm}

\textbf{NIM: \NIM}\\
\textbf{\Author}\\
\vspace{1cm}

{Proposal ini diajukan sebagai usulan pembuatan tugas akhir pada\\
Program Studi Sarjana \Prodi\\
Fakultas Informatika Universitas Telkom}

\vspace{0.5cm}

{Bandung, \Tanggal\quad \Bulan \quad \Date}\\
{Menyetujui}

\vspace{0.5cm}

\begin{center}
Calon Pembimbing
\end{center}

\begin{center}
\vspace{2cm}

\underline{\PembimbingSatu} \\ 
NIP: \NIPPembimbingSatu

\end{center}

\vspace{0.5cm}

}
\pagebreak
\end{titlepage}
\phantomsection
\addcontentsline{toc}{chapter}{Abstrak}
%\pdfbookmark{\abstractname}{Abstrak Indonesia}
\chapter*{Abstrak}

Deteksi sarkasme masih menjadi salah satu tantangan dalam bidang \textit{Natural Language Processing} (NLP) karena makna sarkastis sering kali tidak disampaikan secara langsung dan sangat bergantung pada konteks percakapan. Pada media sosial seperti Reddit, penggunaan bahasa yang informal, singkatan, serta hubungan antar komentar membuat sarkasme semakin sulit dikenali hanya dari makna literal suatu teks. Kondisi tersebut membuat penggunaan model yang mampu memahami konteks menjadi penting, tetapi kebutuhan sumber daya komputasi dari model juga perlu diperhatikan.

Penelitian ini membahas \textit{trade-off} antara performa klasifikasi dan efisiensi komputasi pada dua model transformer, yaitu RoBERTa dan DistilBERT, untuk tugas deteksi sarkasme. Dataset yang digunakan adalah \textit{Self-Annotated Reddit Corpus} (SARC), dengan struktur input berupa pasangan \textit{parent comment} dan \textit{target comment}. Kedua model diuji dalam kondisi eksperimen yang terkontrol tanpa menggunakan \textit{context summarization} maupun metadata tambahan, sehingga perbandingan dapat lebih difokuskan pada karakteristik dari masing-masing model.

Performa klasifikasi dievaluasi menggunakan \textit{accuracy}, \textit{precision}, \textit{recall}, dan \textit{F1-score}. Sementara itu, efisiensi komputasi diukur berdasarkan \textit{training time}, \textit{inference latency}, penggunaan memori GPU, dan jumlah parameter model. Analisis dilakukan untuk melihat seberapa besar perubahan performa yang diperoleh dibandingkan dengan kebutuhan komputasi dari masing-masing model. Penelitian ini tidak bertujuan untuk menentukan satu model yang paling baik secara mutlak, tetapi untuk melihat kondisi penggunaan yang lebih sesuai bagi RoBERTa maupun DistilBERT berdasarkan kebutuhan performa dan keterbatasan sumber daya komputasi.

\vspace{0.5 cm}

\begin{flushleft}
{\textbf{Kata Kunci:} deteksi sarkasme, transformer, \textit{RoBERTa}, \textit{DistilBERT}, \textit{trade-off}, efisiensi komputasi.}
\end{flushleft}
\cleardoublepage
\phantomsection
\addcontentsline{toc}{chapter}{Daftar Isi}
%\pdfbookmark{\contentsname}{Contents}
\tableofcontents

%%%%%%%%%%%%%%%%%%%%%%%%%%%%%%%%%%%%%%%%%%%%%%%%%
% Bab bab penulisan
%%%%%%%%%%%%%%%%%%%%%%%%%%%%%%%%%%%%%%%%%%%%%%%%
\cleardoublepage
\pagenumbering{arabic}
\chapter{Pendahuluan}

\section{Latar Belakang}

Di era digital saat ini, media sosial dan forum diskusi daring telah menjadi sarana komunikasi berbasis teks utama yang digunakan masyarakat untuk menyampaikan opini, berbagi informasi, dan berinteraksi satu sama lain di seluruh dunia. Reddit dan platform serupa lainnya menghasilkan jutaan komentar per hari tentang berbagai topik seperti hiburan, politik, teknologi, isu sosial, dan lain-lain. Besarnya volume komunikasi daring ini telah membuka jalur baru untuk penelitian di bidang Pemrosesan Bahasa Alami (NLP), terutama pada analisis opini dan pemahaman bahasa alami. Namun, beragamnya gaya bahasa pengguna di media sosial juga menghadirkan tantangan linguistik yang cukup kompleks, terutama karena komunikasi daring sering kali mengandung bahasa informal, singkatan, ironi, serta penggunaan makna implisit yang sulit dipahami secara langsung oleh sistem komputasi.

Salah satu jenis bahasa figuratif yang paling umum dalam komunikasi daring adalah sarkasme, yang biasanya beroperasi melalui makna implisit dan kontradiktif. Orang sering menggunakan ekspresi positif atau netral dengan cara yang sarkastis untuk menyampaikan niat negatif, kritik, cemoohan, atau ketidakpuasan. Namun, berbeda dengan bahasa literal, makna yang dimaksud dari sarkasme seringkali berbeda dari makna permukaan kalimat dan oleh karena itu sangat bergantung pada interpretasi kontekstual. Pada media sosial, sarkasme sering kali bergantung pada konteks percakapan, nada komunikasi, dan interaksi sebelumnya antar pengguna. Dengan demikian, sarkasme tidak dapat disimpulkan dengan benar hanya dengan menganalisis kalimat-kalimat individual, terutama di platform diskusi seperti Reddit yang menggunakan struktur percakapan parent-child dengan ketergantungan kontekstual yang panjang \cite{dubey2025unpacking}.

Deteksi sarkasme otomatis merupakan salah satu tantangan utama dalam analisis sentimen dan Pemrosesan Bahasa Alami (NLP) modern. Sarkasme dapat sepenuhnya mengubah sentimen sebenarnya dari suatu kalimat. Model komputasi dapat menghasilkan interpretasi yang salah ketika makna kontekstual tidak dipahami dengan benar. Akibatnya, kegagalan dalam mendeteksi sarkasme dapat menurunkan keandalan berbagai aplikasi seperti opinion mining, analisis umpan balik pelanggan, serta sistem pemantauan media sosial. Selain itu, sifat teks media sosial yang informal dan penuh dengan noise, seperti penggunaan bahasa gaul, singkatan, inkonsistensi tata bahasa, dan ekspresi ambigu, membuat tugas deteksi sarkasme menjadi lebih kompleks. Hal tersebut menunjukkan bahwa deteksi sarkasme memerlukan pendekatan pemahaman bahasa yang mampu menangkap hubungan kontekstual dan informasi semantik implisit dengan lebih baik daripada metode klasifikasi teks konvensional.

Penelitian mengenai deteksi sarkasme telah beralih dari pendekatan berbasis aturan (\textit{rule-based}) dan \textit{traditional machine learning} menuju arsitektur \textit{deep learning} dan berbasis transformer untuk mengatasi keterbatasan pendekatan tradisional. Pendekatan awal sebagian besar bergantung pada pola leksikal, aturan sintaksis, atau fitur yang direkayasa secara manual (\textit{hand-crafted features}). Namun, metode-metode tersebut umumnya tidak berhasil menangkap makna berbasis konteks dan hubungan semantik implisit dari ekspresi sarkastis. Selanjutnya, model pembelajaran mendalam seperti CNN, RNN, dan LSTM diperkenalkan untuk meningkatkan pemahaman kontekstual dengan memanfaatkan mekanisme pembelajaran sekuensial. Namun, perkembangan terkini model berbasis transformer telah merevolusi bidang Pemrosesan Bahasa Alami (NLP) secara signifikan berkat kemampuannya untuk mempelajari representasi kontekstual secara lebih efektif melalui mekanisme \textit{self-attention} \cite{zhang2025enhancing}. Arsitektur berbasis transformer mampu menangkap hubungan semantik antar kata dalam satu kalimat secara keseluruhan, sehingga model mampu memahami sarkasme implisit dan ironi kontekstual dengan lebih baik dalam percakapan media sosial yang tidak terstruktur.

RoBERTa dan DistilBERT merupakan dua model berbasis transformer yang memiliki karakteristik berbeda dalam tugas deteksi sarkasme. RoBERTa dikenal memiliki kemampuan representasi kontekstual yang kuat dan mampu menghasilkan performa klasifikasi yang tinggi. Namun, kemampuan tersebut juga membutuhkan sumber daya komputasi yang relatif lebih besar, seperti penggunaan memori GPU, waktu pelatihan, dan waktu inferensi. Di sisi lain, DistilBERT dikembangkan menggunakan pendekatan *knowledge distillation* untuk menghasilkan model transformer yang lebih ringan dengan tetap mempertahankan kemampuan pemahaman konteks yang kompetitif. Karakteristik tersebut membuat DistilBERT memiliki kebutuhan komputasi yang lebih rendah dan proses inferensi yang lebih cepat \cite{dubey2025unpacking}. Perbedaan ini menunjukkan bahwa pemilihan model tidak selalu dapat ditentukan hanya berdasarkan performa klasifikasi tertinggi, tetapi juga perlu mempertimbangkan kebutuhan komputasi dan kondisi penggunaan model.

Penelitian sebelumnya telah membandingkan RoBERTa dan DistilBERT dari sisi performa klasifikasi maupun efisiensi komputasi. Namun, beberapa penelitian masih menggunakan komponen tambahan seperti *context summarization* dan integrasi metadata dalam proses deteksi sarkasme, sehingga hasil yang diperoleh tidak hanya dipengaruhi oleh karakteristik dari kedua model transformer tersebut. Oleh karena itu, penelitian ini melakukan perbandingan RoBERTa dan DistilBERT pada dataset Self-Annotated Reddit Corpus (SARC) dalam skenario eksperimen yang lebih terkontrol. Analisis yang dilakukan tidak hanya bertujuan untuk melihat model mana yang menghasilkan performa lebih tinggi atau membutuhkan sumber daya lebih sedikit, tetapi juga untuk melihat seberapa besar perubahan performa yang terjadi dibandingkan dengan keuntungan efisiensi komputasi yang diperoleh. Dengan demikian, hasil penelitian diharapkan dapat memberikan gambaran mengenai kondisi penggunaan di mana RoBERTa maupun DistilBERT lebih sesuai berdasarkan kebutuhan performa dan sumber daya komputasi.


\section{Perumusan Masalah}

Berdasarkan latar belakang yang telah dijelaskan sebelumnya, penelitian ini berfokus pada analisis \textit{trade-off} antara performa klasifikasi dan efisiensi komputasi dari model RoBERTa dan DistilBERT untuk tugas deteksi sarkasme. Oleh karena itu, rumusan masalah dalam penelitian ini adalah sebagai berikut:

\begin{enumerate}
    \item Bagaimana perbandingan performa klasifikasi RoBERTa dan DistilBERT dalam mendeteksi sarkasme menggunakan dataset \textit{Self-Annotated Reddit Corpus} (SARC)?
    
    \item Bagaimana perbedaan efisiensi komputasi RoBERTa dan DistilBERT dilihat dari \textit{training time}, \textit{inference latency}, penggunaan memori GPU, dan jumlah parameter model?
    
    \item Bagaimana \textit{trade-off} antara performa klasifikasi dan efisiensi komputasi dari kedua model, serta pada kondisi seperti apa RoBERTa atau DistilBERT lebih sesuai untuk digunakan?
\end{enumerate}

\section{Tujuan}

Penelitian ini bertujuan untuk menganalisis \textit{trade-off} antara performa klasifikasi dan efisiensi komputasi pada model RoBERTa dan DistilBERT untuk tugas deteksi sarkasme menggunakan dataset \textit{Self-Annotated Reddit Corpus} (SARC). Untuk mencapai tujuan tersebut, penelitian ini memiliki beberapa sasaran sebagai berikut:

\begin{enumerate}
    \item Membandingkan performa klasifikasi RoBERTa dan DistilBERT dalam mendeteksi sarkasme pada dataset \textit{Self-Annotated Reddit Corpus} (SARC).
    
    \item Membandingkan efisiensi komputasi RoBERTa dan DistilBERT dalam proses pelatihan dan inferensi.
    
    \item Menganalisis \textit{trade-off} antara performa klasifikasi dan efisiensi komputasi untuk melihat kondisi penggunaan yang lebih sesuai bagi masing-masing model.
\end{enumerate}

\section{Batasan Masalah}

Penelitian ini memiliki beberapa batasan agar pembahasan tetap terfokus pada ruang lingkup yang telah ditentukan, yaitu:

\begin{enumerate}
    \item Penelitian hanya membahas deteksi sarkasme berbasis teks tanpa menggunakan informasi multimodal seperti gambar, audio, maupun video.
    \item Dataset yang digunakan adalah \textit{Self-Annotated Reddit Corpus} (SARC), dengan sekitar 200.000 pasangan \textit{parent comment} dan \textit{target comment} yang digunakan dalam eksperimen.
    \item Data yang digunakan hanya berupa komentar berbahasa Inggris sesuai dengan karakteristik dataset SARC.
    \item Model yang digunakan dalam penelitian ini dibatasi pada RoBERTa dan DistilBERT.
    \item Input model menggunakan \textit{parent comment} dan \textit{target comment} tanpa menggunakan \textit{context summarization}, metadata tambahan, maupun fitur khusus berbasis emoji.
    \item Penelitian berfokus pada analisis \textit{trade-off} antara performa klasifikasi dan efisiensi komputasi, tanpa melakukan implementasi langsung pada lingkungan produksi.
    \item Performa klasifikasi diukur menggunakan accuracy, precision, recall, dan F1-score. Sementara itu, efisiensi komputasi diukur berdasarkan \textit{training time}, \textit{inference latency}, penggunaan memori GPU, dan jumlah parameter model.
    \item Penelitian ini tidak membahas interpretabilitas model maupun pendekatan \textit{Explainable Artificial Intelligence} (XAI).
    \item Proses \textit{fine-tuning} dilakukan dengan pengujian hyperparameter secara terbatas dan tidak mencakup \textit{hyperparameter optimization} secara ekstensif.
\end{enumerate}


\section{Manfaat Penelitian}

Penelitian ini diharapkan dapat memberikan kontribusi secara teoretis dan praktis pada pengembangan sistem deteksi sarkasme berbasis transformer dalam bidang \textit{Natural Language Processing} (NLP). Selain itu, penelitian ini juga diharapkan dapat menawarkan pemahaman yang lebih menyeluruh mengenai analisis \textit{trade-off} antara akurasi klasifikasi dan efisiensi komputasi dalam penerapan model transformer pada konteks media sosial.

\subsection{Manfaat Teoritis}

\begin{enumerate}
    \item Hasil dari penelitian ini diharapkan dapat berkontribusi pada pengembangan penelitian ilmiah terkait deteksi sarkasme berbasis transformer, khususnya dalam hal analisis \textit{trade-off} antara kinerja klasifikasi dan efisiensi komputasi. 
    \item Penelitian ini diharapkan dapat memberikan wawasan akademis yang lebih mendalam mengenai karakteristik arsitektur transformer berukuran besar (\textit{heavyweight}) dan berbobot ringan (\textit{lightweight}), yaitu RoBERTa dan DistilBERT, dalam tugas deteksi sarkasme kontekstual.
    \item Hasil dari penelitian ini diharapkan sebagai tolak ukur bagi penelitian selanjutnya di bidang \textit{Natural Language Processing} yang efisien secara komputasi dan pemahaman bahasa kontekstual pada platform media sosial.
\end{enumerate}


\subsection{Manfaat Praktis}

\begin{enumerate}
    \item Penelitian ini diharapkan dapat memberikan wawasan praktis mengenai implementasi model berbasis transformer untuk deteksi sarkasme pada lingkungan media sosial seperti Reddit.
    \item Hasil dari penelitian ini diharapkan dapat membantu peneliti dan pengembang sistem memahami \textit{trade-off} antara kinerja klasifikasi dan efisiensi komputasi dalam pemilihan model transformer untuk aplikasi NLP di dunia nyata.
    \item Studi ini diharapkan dapat berkontribusi pada pengembangan sistem deteksi sarkasme yang efisien secara komputasi untuk lingkungan analisis media sosial skala besar dan kebutuhan real-time.
\end{enumerate}


\section{Hipotesis}

Tinjauan literatur dan metodologi yang diusulkan untuk penelitian ini mengarah pada perumusan beberapa hipotesis yang akan diuji selama proses eksperimen studi ini. Hipotesis-hipotesis ini dirumuskan berdasarkan karakteristik arsitektur transformer, kemampuan representasi kontekstual model, dan pertimbangan efisiensi komputasi dalam tugas deteksi sarkasme. Hipotesis penelitian ini ditetapkan sebagai berikut:

\begin{enumerate}
    \item Model RoBERTa diperkirakan mencapai kinerja klasifikasi yang lebih tinggi dalam tugas deteksi sarkasme berkat kapabilitas representasi kontekstual yang lebih kuat dan arsitektur transformer yang lebih dalam.
    \item Model DistilBERT diperkirakan memberikan efisiensi komputasi yang lebih baik, seperti waktu inferensi yang lebih cepat dan penggunaan memori GPU yang lebih rendah, karena model ini memanfaatkan arsitektur transformer berbobot ringan berdasarkan \textit{knowledge distillation}.
    \item Meskipun DistilBERT memiliki arsitektur model yang lebih kecil, model ini diperkirakan dapat mempertahankan kinerja klasifikasi yang kompetitif dibandingkan dengan RoBERTa pada dataset \textit{Self-Annotated Reddit Corpus}  (SARC).
    \item Terdapat \textit{trade-off} antara kinerja klasifikasi dan efisiensi komputasi dalam penggunaan RoBERTa dan DistilBERT untuk deteksi sarkasme berbasis transformer.
\end{enumerate}


\section{Rencana Kegiatan}

Rencana kegiatan dalam penelitian ini terdiri dari beberapa tahapan sistematis yang dilakukan untuk mendukung implementasi dan evaluasi model deteksi sarkasme berbasis transformer. Tahap pertama mencakup tinjauan pustaka mengenai deteksi sarkasme, arsitektur transformer, efisiensi komputasi, dan pemahaman bahasa dalam bidang \textit{Natural Language Processing (NLP)}. Setelah itu, dilakukan proses pengumpulan dan persiapan dataset \textit{Self-Annotated Reddit Corpus}  (SARC) sebagai sumber data utama dalam penelitian. Dataset yang telah dikumpulkan kemudian melalui tahap preprocessing data, seperti pembersihan teks, tokenisasi, dan persiapan struktur percakapan kontekstual untuk mendukung tugas deteksi sarkasme.

Langkah selanjutnya mencakup implementasi dan proses fine-tunning model RoBERTa dan DistilBERT untuk klasifikasi sarkasme. Setelah proses implementasi, kedua model akan dievaluasi menggunakan metrik performa klasifikasi dan pengukuran efisiensi komputasi. Selanjutnya, analisis terhadap hasil eksperimen akan dilakukan untuk mengkaji \textit{trade-off} antara kinerja klasifikasi dan efisiensi komputasi dari setiap model transformer. Temuan penelitian dan analisis eksperimental secara keseluruhan akan dimuat dalam laporan penelitian akhir.


\section{Jadwal Kegiatan}

Jadwal pelaksanaan penelitian disusun berdasarkan tahapan metodologi penelitian yang telah dirancang. Adapun jadwal kegiatan penelitian yang akan dilakukan ditunjukkan pada Tabel 1.1.

\begin{table}[h!]
\centering
\caption{Jadwal kegiatan proposal tugas akhir}
\label{tab:jadwal_kegiatan}

\begin{tabular}{|c|m{3.5cm}|*{24}{c|}}
\hline

\multirow{2}{*}{\textbf{No}} &
\multirow{2}{*}{\textbf{Kegiatan}} &
\multicolumn{24}{c|}{\textbf{Bulan ke-}} \\

\cline{3-26}

& &
\multicolumn{4}{c|}{\textbf{1}} &
\multicolumn{4}{c|}{\textbf{2}} &
\multicolumn{4}{c|}{\textbf{3}} &
\multicolumn{4}{c|}{\textbf{4}} &
\multicolumn{4}{c|}{\textbf{5}} &
\multicolumn{4}{c|}{\textbf{6}} \\

\hline

1 &
Studi Literatur &
\cellcolor{blue!25} & \cellcolor{blue!25} & \cellcolor{blue!25} & \cellcolor{blue!25} &
\cellcolor{blue!25} & \cellcolor{blue!25} & \cellcolor{blue!25} & \cellcolor{blue!25} &
{} & {} & {} & {} &
{} & {} & {} & {} &
{} & {} & {} & {} &
{} & {} & {} & {} \\
\hline

2 &
Pengumpulan Dataset &
\cellcolor{blue!25} & \cellcolor{blue!25} & \cellcolor{blue!25} & \cellcolor{blue!25} &
\cellcolor{blue!25} & \cellcolor{blue!25} & \cellcolor{blue!25} & \cellcolor{blue!25} &
{} & {} & {} & {} &
{} & {} & {} & {} &
{} & {} & {} & {} &
{} & {} & {} & {} \\
\hline

3 &
Analisis dan Perancangan Sistem &
{} & {} & {} & {} &
\cellcolor{blue!25} & \cellcolor{blue!25} & \cellcolor{blue!25} & \cellcolor{blue!25} &
\cellcolor{blue!25} & \cellcolor{blue!25} & \cellcolor{blue!25} & \cellcolor{blue!25} &
{} & {} & {} & {} &
{} & {} & {} & {} &
{} & {} & {} & {} \\
\hline

4 &
Implementasi dan \textit{Fine-Tuning} Model &
{} & {} & {} & {} &
{} & {} & {} & {} &
\cellcolor{blue!25} & \cellcolor{blue!25} & \cellcolor{blue!25} & \cellcolor{blue!25} &
\cellcolor{blue!25} & \cellcolor{blue!25} & \cellcolor{blue!25} & \cellcolor{blue!25} &
{} & {} & {} & {} &
{} & {} & {} & {} \\
\hline

5 &
Evaluasi dan Analisis Hasil &
{} & {} & {} & {} &
{} & {} & {} & {} &
{} & {} & {} & {} &
\cellcolor{blue!25} & \cellcolor{blue!25} & \cellcolor{blue!25} & \cellcolor{blue!25} &
\cellcolor{blue!25} & \cellcolor{blue!25} & \cellcolor{blue!25} & \cellcolor{blue!25} &
{} & {} & {} & {} \\
\hline

6 &
Penyusunan Laporan &
{} & {} & {} & {} &
\cellcolor{blue!25} & \cellcolor{blue!25} & \cellcolor{blue!25} & \cellcolor{blue!25} &
\cellcolor{blue!25} & \cellcolor{blue!25} & \cellcolor{blue!25} & \cellcolor{blue!25} &
\cellcolor{blue!25} & \cellcolor{blue!25} & \cellcolor{blue!25} & \cellcolor{blue!25} &
\cellcolor{blue!25} & \cellcolor{blue!25} & \cellcolor{blue!25} & \cellcolor{blue!25} &
\cellcolor{blue!25} & \cellcolor{blue!25} & \cellcolor{blue!25} & \cellcolor{blue!25} \\
\hline

\end{tabular}

\end{table}

\newpage
%
\chapter{Kajian Pustaka}

\section{Konsep dan Karakteristik Sarkasme dalam Teks}

Sarkasme merupakan salah satu bentuk bahasa figuratif yang cukup unik karena makna harfiahnya sering kali berlawanan dengan maksud yang ingin disampaikan. Dalam komunikasi sehari-hari, sarkasme lazim digunakan sebagai cara tidak langsung untuk menyampaikan kritik, ejekan, atau ketidakpuasan, sehingga pesan yang disampaikan menjadi tidak mudah ditangkap secara langsung. Sifat ambigu inilah yang kemudian menjadikan sarkasme sebagai salah satu fenomena paling menantang dalam pemrosesan bahasa alami (NLP), terutama dalam konteks analisis sentimen.

Dalam teks tertulis, sarkasme dapat dikenali melalui beberapa ciri yang membedakannya dari bentuk komunikasi lain. Ciri yang paling mendasar adalah adanya pertentangan makna atau inkonsistensi pragmatis. Penutur secara sengaja menyimpang dari komunikasi normal yang berlaku, misalnya dengan memilih diksi positif untuk menyampaikan maksud yang justru negatif. Di platform media sosial, bahasa sarkasme juga kerap hadir dalam bentuk tidak baku, mulai dari penggunaan slang, singkatan, hingga ekspresi. Selain itu, penanda visual seperti tanda baca berlebihan, huruf kapital, dan pengulangan huruf tertentu sering muncul sebagai upaya penutur mempertegas emosi yang ingin disampaikan. Emoji pun tidak jarang dimanfaatkan sebagai petunjuk tambahan untuk menandai nada sarkastis, meski sifatnya yang polisemi justru kerap menimbulkan ambiguitas \cite{grover2026emoji}.

Dibandingkan tugas klasifikasi teks pada umumnya, deteksi sarkasme jelas jauh lebih kompleks. Pendekatan analisis sentimen yang konvensional umumnya mengandalkan kata-kata bermuatan emosional untuk menentukan polaritas suatu teks. Pendekatan tersebut menjadi kurang efektif ketika diterapkan pada teks sarkastis, sebab makna yang sesungguhnya sering kali tidak tercermin dari permukaan teks. Mendeteksi sarkasme membutuhkan pemahaman konteks yang lebih luas, mencakup hubungan antar kalimat, konteks percakapan sebelumnya, hingga informasi eksternal yang relevan.

Dalam perspektif komputasi, tantangan yang dihadapi bahkan jauh lebih berlapis. Tidak tersedianya sinyal non-verbal seperti nada suara dan ekspresi wajah menjadi kendala yang cukup fundamental, mengingat dalam komunikasi tatap muka kedua elemen tersebut berperan besar dalam mengidentifikasi sarkasme. Di luar itu, menafsirkan sarkasme memang memerlukan kepekaan terhadap nuansa makna yang tidak selalu mudah ditangkap, bahkan oleh manusia sekalipun. Tantangan tersebut semakin besar karena data media sosial umumnya tidak terstruktur, penuh derau, dan mengalami ketidakseimbangan kelas antara sampel sarkastis dan non-sarkastis \cite{dubey2025unpacking}. Tidak hanya itu, sarkasme sering kali hadir bersama konten multimodal seperti gambar, video, atau emoji, sementara banyak pendekatan yang ada masih terbatas pada analisis teks saja \cite{dhumpati2025sarcasm}. Di sisi lain, model-model canggih yang dirancang untuk menangkap relasi kontekstual secara mendalam, seperti model berbasis transformer dan jaringan saraf dua arah, umumnya memerlukan kapasitas komputasi yang besar, yang menjadi hambatan nyata dalam penerapannya di dunia nyata.

Berdasarkan berbagai karakteristik dan tantangan tersebut, deteksi sarkasme pada dasarnya merupakan permasalahan yang tidak bisa diselesaikan dengan pendekatan konvensional yang sederhana. Dibutuhkan model yang mampu menangkap ketergantungan kontekstual secara mendalam, namun tetap mempertimbangkan efisiensi komputasi agar layak dan optimal diterapkan dalam konteks praktis atau \textit{real-time} \cite{dubey2025unpacking}.

\section{Pendekatan dalam Deteksi Sarkasme}

Metode deteksi sarkasme telah berkembang secara signifikan seiring dengan meningkatnya kompleksitas data tekstual di media sosial dan kebutuhan yang semakin besar bagi model untuk memahami informasi kontekstual secara lebih efektif. Metode-metode ini berkisar dari pendekatan berbasis aturan yang sederhana hingga arsitektur berbasis transformer yang canggih yang dapat menangkap hubungan semantik dalam konteks. Setiap pendekatan memiliki karakteristik, kekuatan, dan kelemahan yang berbeda dalam menangani tugas deteksi sarkasme, khususnya terkait \textit{trade-off} antara akurasi analitis dan efisiensi sumber daya komputasi.

\subsection{Pendekatan Berbasis (\textit{rule-based}) dan \textit{Machine Learning}}

Pendekatan awal untuk deteksi sarkasme sebagian besar didasarkan pada pendekatan berbasis aturan dan statistik yang sangat bergantung pada fitur-fitur yang dirancang secara manual. Metode-metode ini melibatkan penggunaan kamus leksikal, TF-IDF, n-gram, dan penandaan Part-of-Speech (POS) untuk mengidentifikasi pola linguistik dalam data tekstual. Selain itu, beberapa penelitian telah menggunakan metode deteksi kontradiksi dan ketidakkonsistenan sentimen untuk menemukan konflik antara ekspresi positif dan situasi kontekstual negatif dalam sebuah kalimat.

Pendekatan berbasis (\textit{rule-based}) cukup efektif untuk mengenali pola sarkasme yang eksplisit, namun performanya terbatas karena aturan-aturan tersebut dirancang secara manual dan bersifat kaku. Terlebih lagi, teknik-teknik ini gagal beradaptasi dengan variasi bahasa yang dinamis dan untuk memahami makna kontekstual implisit, yang biasanya merupakan inti dari ekspresi sarkastis. Keterbatasan ini menyebabkan meluasnya penggunaan algoritma pembelajaran mesin tradisional dalam tugas-tugas deteksi sarkasme. Beberapa algoritma yang paling populer adalah Support Vector Machine (SVM), Naive Bayes, K-Nearest Neighbors (KNN), Regresi Logistik, dan algoritma berbasis ensemble seperti Random Forest, XGBoost, dan AdaBoost. Model-model ini digunakan untuk tugas klasifikasi berdasarkan fitur-fitur linguistik yang telah diekstraksi sebelumnya.

Pendekatan berbasis aturan cukup efektif untuk mengenali pola sarkasme yang eksplisit, namun performanya terbatas karena aturan-aturan tersebut dirancang secara manual dan bersifat kaku. Selain itu, pendekatan ini memerlukan proses perancangan, pemeliharaan, dan pembaruan aturan secara terus-menerus agar dapat mengikuti perkembangan variasi bahasa yang dinamis. Keterbatasan tersebut menyebabkan pendekatan berbasis aturan kurang mampu menangkap makna kontekstual dan hubungan semantik implisit yang sering menjadi karakteristik utama ekspresi sarkastis. Untuk mengatasi keterbatasan tersebut, berbagai penelitian mulai beralih ke pendekatan machine learning yang mampu mempelajari pola bahasa secara otomatis dari data. Namun, model pembelajaran mesin tradisional masih menghadapi tantangan dalam memodelkan ketergantungan kontekstual dan ambiguitas semantik karena representasi teks yang digunakan umumnya belum mampu menangkap hubungan kontekstual secara mendalam \cite{dubey2025unpacking}.

\subsection{Pendekatan Berbasis \textit{deep learning}}

Perkembangan \textit{deep learning} membantu deteksi sarkasme. Hal ini memungkinkan model mempelajari representasi fitur otomatis dan mengurangi kebutuhan rekayasa fitur manual. Salah satu arsitektur yang banyak digunakan adalah \textit{Convolutional Neural Network} (CNN), yang efektif dalam menangkap pola lokal dan hubungan spasial antar kata dalam sebuah kalimat.

Selain CNN, arsitektur berbasis \textit{Recurrent Neural Network} (RNN) seperti \textit{Long Short-Term Memory} (LSTM) dan \textit{Gated Recurrent Unit} (GRU) juga banyak digunakan untuk deteksi sarkasme. Model-model ini dirancang untuk menangani data sekuensial dan mengingat konteks dari waktu yang lama. Inovasi lebih lanjut menghasilkan versi dua arah seperti \textit{Bidirectional LSTM} (BiLSTM) dan \textit{Bidirectional GRU} (BiGRU), yang memungkinkan model untuk menangkap informasi kontekstual dari urutan sebelumnya dan sesudahnya secara bersamaan.

Selain itu, terdapat beberapa penelitian yang mengusulkan arsitektur \textit{hybrid} yang menggabungkan model CNN dan LSTM untuk menangkap fitur spasial dan temporal secara bersama-sama \cite{prakash2025sarcasm}. Sebagai tambahan, arsitektur \textit{deep learning} telah dilengkapi dengan mekanisme atensi untuk mengidentifikasi kata-kata yang merupakan indikator kuat sarkasme dengan lebih baik. Teknik seperti Attention Mechanism-Based Bidirectional LSTM-GRU (AM-BLSTM-GRU) telah melampaui metode \textit{deep learning} tradisional dalam hal pengetahuan kontekstual \cite{dhumpati2025sarcasm}.

Model \textit{deep learning} telah membuat kemajuan dalam pemahaman semantik dan kontekstual, namun mereka masih memiliki keterbatasan dalam menangani hubungan kontekstual yang sangat kompleks dan dependensi jangka panjang. Selain itu, persyaratan komputasi yang tinggi tetap menjadi tantangan untuk penggunaan praktis \cite{dhumpati2025sarcasm}.

\subsection{Pendekatan Berbasis Transformer}

Perkembangan terbaru dalam deteksi sarkasme didominasi oleh pendekatan transfer learning dan arsitektur berbasis transformer. Berbeda dengan arsitektur rekuren tradisional, transformer memproses semua hubungan antar kata secara paralel di seluruh kalimat, melalui mekanisme \textit{self-attention}. Mekanisme ini memungkinkan model untuk menghasilkan \textit{contextual embedding}, di mana makna suatu kata berubah sesuai dengan konteksnya.

Arsitektur berbasis \textit{transformer} seperti BERT, RoBERTa, dan DistilBERT telah menunjukkan performa superior dalam deteksi sarkasme karena kemampuannya menangkap informasi kontekstual dua arah. Secara spesifik, RoBERTa yang dikembangkan dengan proses pelatihan mendalam mampu menghasilkan akurasi klasifikasi yang tinggi, namun menuntut sumber daya komputasi yang signifikan, seperti konsumsi memori GPU yang besar dan waktu pelatihan yang lama. Sebagai alternatif, DistilBERT yang dirancang melalui proses \textit{knowledge distillation} menawarkan arsitektur yang lebih ringan dan waktu inferensi yang jauh lebih cepat. Meskipun terdapat sedikit penurunan performa prediksinya dibandingkan RoBERTa, efisiensi DistilBERT menjadikannya solusi yang lebih layak untuk penerapan sistem dengan sumber daya terbatas dan kebutuhan \textit{real-time} \cite{dubey2025unpacking}.

Perbedaan karakteristik antara RoBERTa dan DistilBERT menunjukan adanya \textit{trade-off} antara akurasi dan efisiensi komputasi. Oleh karena itu, kedua model tersebut menjadi fokus utama dalam penelitian ini untuk menganalisis keseimbangan antara performa klasifikasi dan kebutuhan komputasi dalam deteksi sarkasme berbasis teks.

\subsection{Pendekatan Unimodal dan Multimodal}

Selain pendekatan berbasis teks (\textit{unimodal}), beberapa penelitian terbaru mulai mengembangkan pendekatan \textit{multimodal} menggunakan informasi tambahan seperti audio, gambar, dan emoji untuk menangkap konteks non-verbal. Meskipun pendekatan multimodal dapat memperkaya pemahaman kontekstual, integrasi berbagai tipe data ini secara substansial membuat sistem menjadi jauh lebih kompleks dan menuntut sumber daya komputasi yang besar \cite{dhumpati2025sarcasm}. Oleh karena itu, penelitian ini membatasi ruang lingkup pada pendekatan \textit{unimodal} berbasis teks. Pembatasan ini dilakukan agar analisis dapat terfokus secara mendalam pada evaluasi \textit{trade-off} antara akurasi pemahaman konteks dan efisiensi komputasi pada arsitektur transformer, khususnya model DistilBERT dan RoBERTa.

Meskipun pendekatan multimodal dapat meningkatkan pemahaman kontekstual, pendekatan tersebut juga secara substansial membuat sistem menjadi lebih kompleks dan membutuhkan sumber daya komputasi yang lebih besar \cite{dhumpati2025sarcasm}. Oleh karena itu, penelitian ini berfokus pada pendekatan \textit{unimodal} berbasis teks agar analisis dapat difokuskan pada evaluasi \textit{trade-off} antara akurasi dan efisiensi komputasi pada arsitektur transformer, khususnya DistilBERT dan RoBERTa. 

\section{Metrik Evaluasi Kinerja Model}
\subsection{Confusion Matrix}

Dalam penelitian klasifikasi teks, khususnya pada deteksi sarkasme, confusion matrix sering digunakan untuk evaluasi kinerja model. \textit{Confusion matrix} merupakan matriks evaluasi yang digunakan untuk membandingkan hasil prediksi model dengan label aktual dari data yang diamati. Matriks ini memungkinkan analisis yang lebih rinci mengenai kinerja model dengan melihat jumlah prediksi yang benar dan salah untuk setiap kelas.

\textit{Confusion matrix} memiliki empat komponen utama, yaitu \textit{True Positive} (TP), \textit{True Negative} (TN), \textit{False Positive} (FP), \textit{False Negative} (FN). \textit{True Positive} menunjukan jumlah teks sarkastis yang berhasil diklasifikasikan benar sebagai teks sarkasme oleh model. Sebaliknya, \textit{True Negative} menunjukan jumlah teks non-sarkastis yang berhasil diklasifikasikan benar sebagai teks non-sarkastis oleh model. Sementara itu, \textit{False Positive} adalah ketika model salah mengklasifikasikan teks non-sarkastis sebagai sarkastis, dan \textit{False Negative} adalah ketika teks sarkastis tidak terdeteksi dan diklasifikasikan sebagai teks normal.

Dalam konteks deteksi sarkasme, analisis mendetail terhadap nilai TP, TN, FP, dan FN sangat penting karena klasifikasi sering kali terkecoh oleh makna sarkasme yang implisit dan ambigu.

\subsection{Accuracy}
Accuracy atau akurasi merupakan metrik evaluasi yang mengukur proporsi prediksi yang diklasifikasikan dengan benar dari total jumlah data yang diuji. Metrik ini memberikan gambaran umum mengenai kemampuan model dalam melakukan klasifikasi secara keseluruhan.

Rumus accuracy dapat dinyatakan sebagai berikut:

\begin{equation}
\text{Accuracy} = \frac{\text{TP} + \text{TN}}{\text{TP} + \text{TN} + \text{FP} + \text{FN}}
\end{equation}

Semakin tinggi nilai accuracy, semakin banyak data yang diklasifikasikan dengan benar oleh model. Namun, accuracy memiliki keterbatasan ketika berhadapan dengan kumpulan data yang tidak seimbang. Dalam deteksi sarkasme, jumlah data non-sarkastis sering kali lebih banyak daripada teks sarkastis dalam kumpulan data \cite{dhumpati2025sarcasm}. Dengan demikian, sebuah model dapat mencapai akurasi tinggi hanya dengan memprediksi sebagian besar data sebagai kelas mayoritas tanpa benar-benar mempelajari karakteristik ekspresi sarkastis.

\subsection{Precision}
Precision digunakan untuk mengukur seberapa akurat prediksi positif yang dihasilkan oleh model. Dalam konteks deteksi sarkasme, precision menunjukan seberapa banyak teks sarkastis diprediksi benar sebagai sarkastis.

Rumus precision ditunjukkan sebagai berikut:

\begin{equation}
\text{Precision} = \frac{\text{TP}}{\text{TP} + \text{FP}}
\end{equation}

Nilai precision yang tinggi menunjukkan bahwa model dapat mengurangi \textit{false positive}. Ini penting untuk menghindari salah klasifikasi teks non-sarkastis sebagai sarkastis dalam konteks deteksi sarkasme. Dengan demikian, precision merupakan salah satu indikator penting untuk mengevaluasi ketepatan model dalam memprediksi sarkasme.


\subsection{Recall}
Recall merupakan metrik evaluasi yang digunakan untuk mengukur kemampuan model dalam mengidentifikasi semua data positif yang tersedia. Dalam penelitian deteksi sarkasme, recall menunjukan kemampuan model untuk mengidentifikasi teks sarkastis di antara semua teks yang sebenarnya sarkastis. 

Rumus recall ditunjukkan sebagai berikut:

\begin{equation}
\text{Recall} = \frac{\text{TP}}{\text{TP} + \text{FN}}
\end{equation}

Semakin tinggi nilai recall, semakin sedikit false negatives yang dihasilkan oleh model. Recall menjadi indikator penting dalam deteksi sarkasme karena banyak ekspresi sarkastis yang disampaikan secara implisit dan sulit diidentifikasi. Oleh karena itu, model dengan recall yang baik diharapkan dapat menangkap lebih banyak pola sarkasme tersembunyi dalam teks.

\subsection{F1-Score}
F1-score adalah rata-rata harmonik dari precision dan recall. Metrik ini digunakan untuk memberikan evaluasi kinerja model yang lebih seimbang, terutama dalam kasus di mana dataset memiliki distribusi kelas yang tidak seimbang.

Rumus F1-Score ditunjukkan sebagai berikut:

\begin{equation}
\text{F1-Score} = 2 \times \frac{\text{Precision} \times \text{Recall}}{\text{Precision} + \text{Recall}}
\end{equation}

F1-score hanya akan menghasilkan nilai yang tinggi apabila \textit{precision} dan \textit{recall} berada dalam kondisi seimbang. Mengingat kompleksitas deteksi sarkasme sering kali melibatkan tarik-ulur (\textit{trade-off}) antara mempertahankan presisi dan kepekaan menangkap konteks tersembunyi, \textit{F1-score} dioptimalkan sebagai metrik evaluasi utama. Metrik ini memberikan penilaian yang jauh lebih komprehensif dan adil dibandingkan \textit{accuracy}, khususnya dalam menangani dataset sarkasme di dunia nyata yang secara alami tidak seimbang \cite{dubey2025unpacking}.

\subsection{Computational Efficiency Metrics}
Selain mengevaluasi performa klasifikasi, efisiensi komputasi model juga dipelajari dalam penelitian ini di samping kinerja klasifikasi. Efisiensi komputasi merupakan pertimbangan penting karena arsitektur berbasis transformer biasanya memiliki jumlah parameter yang besar dan struktur arsitektur yang kompleks, sehingga membutuhkan biaya komputasi yang tinggi.

Beberapa metrik efisiensi komputasi seperti waktu pelatihan, waktu inferensi, penggunaan memori, dan ukuran parameter digunakan dalam penelitian ini. Waktu pelatihan mengacu pada waktu yang dibutuhkan model untuk melakukan proses pelatihan, dan waktu inferensi mengacu pada latensi untuk membuat prediksi pada data yang tidak terlihat. Penggunaan memori di sini mengacu pada jumlah memori komputasi yang digunakan selama eksekusi. Ukuran parameter mengacu pada kompleksitas dan ukuran arsitektur model.

Evaluasi terhadap efisiensi komputasi menjadi penting dalam penerapan model pada dunia nyata (\textit{real-world environment}), terutama dalam sistem dengan keterbatasan sumber daya. Dalam penelitian ini, RoBERTa diperkirakan mampu mencapai kinerja klasifikasi yang lebih baik karena memiliki kapasitas pemahaman konteks yang lebih besar. Namun, peningkatan ini membutuhkan sumber daya komputasi yang lebih besar. Sementara itu, DistilBERT dirancang untuk menjadi model yang ringan dan efisien dengan lebih sedikit parameter yang dapat memungkinkan inferensi lebih cepat dan penggunaan memori yang lebih rendah, meskipun terdapat kemungkinan penurunan performa klasifikasi dibandingkan dengan RoBERTa \cite{dubey2025unpacking}.

Berdasarkan karakteristik tersebut, studi ini dilakukan untuk menganalisis trade-off antara akurasi dan efisiensi komputasi deteksi sarkasme berbasis transformer, khususnya menggunakan DistilBERT dan RoBERTa.


\section{Penelitian Terdahulu}

Penelitian mengenai deteksi sarkasme telah mendapat perhatian yang cukup besar dalam beberapa tahun terakhir, khususnya seiring meningkatnya penggunaan model berbasis transformer. Sejumlah penelitian sebelumnya telah mengeksplorasi berbagai pendekatan, mulai dari peningkatan performa klasifikasi, pemanfaatan konteks percakapan, hingga efisiensi komputasi model. Ringkasan penelitian-penelitian tersebut yang relevan dengan topik ini disajikan pada Tabel \ref{tab:penelitian_terdahulu}.


\begin{table}[H]
\centering
\caption{Perbandingan Penelitian Terdahulu}
\label{tab:penelitian_terdahulu}
\resizebox{\textwidth}{!}{
\begin{tabular}{|p{2.8cm}|p{2.5cm}|p{2.5cm}|p{3cm}|p{3.2cm}|p{3cm}|p{3cm}|}
\hline
\textbf{Peneliti} &
\textbf{Dataset} &
\textbf{Model} &
\textbf{Fokus Penelitian} &
\textbf{Hasil Utama} &
\textbf{Keterbatasan} &
\textbf{Research Gap} \\
\hline

Jiménez et al. (2026)
&
iSarcasmEval, Sarcasm Corpus V2
&
BERT, RoBERTa, ALBERT, DistilBERT, DeBERTa
&
Perbandingan performa beberapa model berbasis BERT untuk deteksi sarkasme.
&
DeBERTa menunjukkan performa paling baik pada dataset yang digunakan, sementara DistilBERT menunjukkan keunggulan dari sisi efisiensi waktu.
&
Tidak menggunakan dataset SARC dan pembahasan efisiensi komputasi masih terbatas pada aspek tertentu.
&
Masih diperlukan kajian trade-off performa dan efisiensi komputasi pada data percakapan Reddit yang lebih kontekstual.
\\
\hline

Ullah et al. (2026)
&
News Headlines dan The Onion
&
RoBERTa, DistilBERT, ALBERT, Soft Voting Ensemble
&
Deteksi sarkasme pada variasi bahasa formal dan informal serta interpretabilitas model menggunakan \textit{LIME}.
&
Pendekatan ensemble menunjukkan performa yang baik, sedangkan DistilBERT tetap kompetitif sebagai model yang lebih ringan.
&
Tidak menggunakan dataset percakapan kontekstual seperti SARC dan belum berfokus pada analisis efisiensi komputasi secara mendalam.
&
Masih diperlukan evaluasi model pada data percakapan media sosial yang lebih kontekstual dengan mempertimbangkan performa dan efisiensi komputasi.
\\
\hline

Ataei et al. (2020)
&
Twitter dan SARC
&
NBSVM, BERT, XLNet, Bi-GRU-CNN, LCF-BERT
&
Integrasi \textit{Aspect-Based Sentiment Analysis} (ABSA) dan model Transformer untuk deteksi sarkasme.
&
Model berbasis transformer menunjukkan performa yang lebih baik dibandingkan beberapa pendekatan tradisional.
&
Fokus utama penelitian berada pada peningkatan performa klasifikasi dan belum mengevaluasi efisiensi komputasi secara khusus.
&
Masih diperlukan evaluasi yang mempertimbangkan kebutuhan sumber daya komputasi dan kelayakan implementasi model pada kondisi praktis.
\\
\hline

Dubey et al. (2026)
&
News Headlines, Mustard, dan SARC
&
RoBERTa, DistilBERT dengan \textit{BART Context Summarization}
&
Analisis performa dan efisiensi arsitektur Transformer untuk deteksi sarkasme.
&
RoBERTa menunjukkan performa klasifikasi yang tinggi, sedangkan DistilBERT menunjukkan efisiensi komputasi yang lebih baik.
&
Menggunakan komponen tambahan seperti \textit{BART} dan metadata sehingga evaluasi efisiensi belum sepenuhnya merepresentasikan perbandingan langsung arsitektur model.
&
Masih diperlukan evaluasi trade-off yang lebih terisolasi antara RoBERTa dan DistilBERT tanpa komponen tambahan.
\\
\hline

Zhang et al. (2025)
&
SARC
&
BERT Multi-Head Attention, GPT-3, Claude-2, Llama-2
&
Optimalisasi Transformer menggunakan \textit{Multi-Head Attention} untuk meningkatkan sensitivitas konteks.
&
Model BERT dengan \textit{Multi-Head Attention} menunjukkan performa yang kompetitif pada dataset SARC dan mampu mengungguli beberapa model LLM.
&
Penelitian lebih berfokus pada peningkatan performa klasifikasi dan belum membahas alternatif model ringan secara mendalam.
&
Masih diperlukan kajian mengenai keseimbangan antara peningkatan performa dan efisiensi komputasi pada model Transformer yang lebih ringan.
\\
\hline

\end{tabular}
}
\end{table}


Berdasarkan Tabel \ref{tab:penelitian_terdahulu}, Sebagian besar penelitian yang ada cenderung berfokus pada pemanfaatan arsitektur berbasis transformer untuk meningkatkan performa deteksi sarkasme. Model-model seperti BERT, RoBERTa, DeBERTa, dan berbagai variannya terbukti mampu menangkap konteks secara lebih baik sehingga menghasilkan akurasi klasifikasi yang cukup tinggi. Tak sedikit pula penelitian yang mencoba mengombinasikan pendekatan tersebut dengan attention mechanism, ensemble learning, maupun penggabungan informasi kontekstual guna memperkuat kemampuan deteksi.

Namun demikian, fokus utama pada sebagian besar penelitian tersebut masih berkutat pada pencapaian performa klasifikasi semata. Aspek efisiensi komputasi seperti waktu pelatihan, latensi inferensi, konsumsi memori GPU, serta ukuran model justru jarang dibahas secara mendalam. Beberapa penelitian yang sudah menyinggung efisiensi komputasi pun umumnya melibatkan komponen tambahan seperti context summarization atau fitur eksternal, sehingga gambaran mengenai trade-off antara performa dan efisiensi pada arsitektur transformer itu sendiri masih belum terlalu jelas.

\section{Research Gap dan Posisi Penelitian}

Penelitian sebelumnya telah membandingkan RoBERTa dan DistilBERT untuk deteksi sarkasme, termasuk pada dataset SARC. Hasil penelitian menunjukkan bahwa RoBERTa cenderung memiliki performa klasifikasi yang lebih tinggi, sedangkan DistilBERT lebih efisien dari sisi komputasi. Namun, beberapa penelitian juga menggunakan komponen tambahan seperti context summarization dan metadata, sehingga hasil yang diperoleh tidak hanya dipengaruhi oleh kedua model tersebut.

Penelitian ini melakukan perbandingan RoBERTa dan DistilBERT pada kondisi eksperimen yang lebih terkontrol menggunakan input parent comment dan target comment tanpa context summarization maupun metadata tambahan. Evaluasi dilakukan berdasarkan performa klasifikasi serta kebutuhan komputasi dari masing-masing model.

Penelitian ini tidak bertujuan menentukan satu model yang terbaik secara mutlak, tetapi melihat seberapa besar performa yang diperoleh dibandingkan dengan sumber daya komputasi yang dibutuhkan. Hasilnya diharapkan dapat memberikan gambaran kapan RoBERTa lebih sesuai digunakan dan kapan DistilBERT menjadi pilihan yang lebih efisien.

%
\chapter{Metodologi dan Desain Sistem}

\section{Perancangan Alur Sistem}

\begin{figure}[H]
\centering

\tikzstyle{startstop} = [
rectangle,
rounded corners,
minimum width=2.8cm,
minimum height=0.8cm,
text centered,
draw=black,
fill=red!20
]

\tikzstyle{process} = [
rectangle,
minimum width=3.5cm,
minimum height=0.8cm,
text centered,
draw=black,
fill=orange!20
]

\tikzstyle{arrow} = [
thick,
->,
>=stealth
]

\begin{tikzpicture}[
scale=0.85,
transform shape,
node distance=1.1cm
]

\node (start) [startstop] {Start};

\node (dataset) [process, below of=start]
{Dataset Collection};

\node (preprocess) [process, below of=dataset]
{Data Preprocessing};

\node (split) [process, below of=preprocess]
{\shortstack{Data Splitting\\\small{(Train / Validation / Test)}}};

\node (roberta) [process,
below left of=split,
xshift=-1.8cm,
yshift=-0.8cm]
{RoBERTa Fine-Tuning};

\node (distilbert) [process,
below right of=split,
xshift=1.8cm,
yshift=-0.8cm]
{DistilBERT Fine-Tuning};

\node (evaluation) [process,
below of=split,
yshift=-2.8cm]
{Performance Evaluation};

\node (efficiency) [process, below of=evaluation]
{Computational Efficiency Measurement};

\node (tradeoff) [process, below of=efficiency]
{Trade-off Analysis};

\node (end) [startstop, below of=tradeoff]
{Conclusion};

\draw [arrow] (start) -- (dataset);

\draw [arrow] (dataset) -- (preprocess);

\draw [arrow] (preprocess) -- (split);

\draw [arrow] (split) -- (roberta);

\draw [arrow] (split) -- (distilbert);

\draw [arrow] (roberta) |- (evaluation);

\draw [arrow] (distilbert) |- (evaluation);

\draw [arrow] (evaluation) -- (efficiency);

\draw [arrow] (efficiency) -- (tradeoff);

\draw [arrow] (tradeoff) -- (end);

\end{tikzpicture}

\caption{Research flow of the proposed system}
\label{fig:researchflow}

\end{figure}

Studi ini dilakukan untuk menganalisis \textit{trade-off} antara akurasi klasifikasi dan efisiensi komputasi dalam tugas deteksi sarkasme menggunakan dua model berbasis transformer, yaitu RoBERTa dan DistilBERT. Penelitian ini berfokus pada evaluasi efektivitas kedua model tersebut dalam memahami konteks sarkastis pada teks media sosial. Selain itu, kebutuhan komputasi untuk pelatihan dan inferensi juga dipertimbangkan. Pendekatan ini dipilih karena model arsitektur transformer modern umumnya menunjukkan performa klasifikasi yang tinggi tetapi mahal secara komputasi sehingga kurang cocok untuk implementasi secara \textit{\textit{real-time}} \cite{dubey2025unpacking}.

Alur kerja penelitian umumnya dimulai dengan mengumpulkan dataset sarkasme dari dataset benchmark publik dan media sosial. Data yang terkumpul kemudian diproses melalui serangkaian tahap preprocessing, termasuk penghapusan data \textit{noise}, normalisasi teks, tokenisasi, dan persiapan input agar kompatibel dengan model berbasis transformer. Setelah langkah preprocessing, dataset dibagi menjadi data pelatihan, validasi, dan pengujian. Selanjutnya, pelatihan dilakukan menggunakan skenario yang sama untuk menyempurnakan model RoBERTa dan DistilBERT secara terpisah guna memberikan proses evaluasi yang adil dan konsisten.

Tahap selanjutnya adalah mengevaluasi performa kedua model menggunakan beberapa metrik klasifikasi yang berbeda seperti accuracy, precision, recall, dan F1-score. Studi ini juga mengukur efisiensi komputasi model dalam hal waktu pelatihan, waktu inferensi, penggunaan memori, dan jumlah parameter dalam model. Hasil yang diperoleh dari evaluasi performa dan pengukuran efisiensi komputasi dari kedua model dianalisis untuk melihat hubungan antara kualitas klasifikasi dan kebutuhan sumber daya komputasi untuk setiap model. Gambaran umum tahapan penelitian secara keseluruhan diilustrasikan pada Gambar 3.1. Deskripsi masing-masing tahapan diuraikan pada subbab berikut.

\section{Dataset dan Sumber Data}

Dataset utama yang digunakan dalam penelitian ini adalah \textit{Self-Annotated Reddit Corpus} (SARC). SARC merupakan salah satu dataset \textit{benchmark} yang banyak digunakan dalam deteksi sarkasme karena diekstrak langsung dari forum diskusi Reddit, mencakup komunitas seperti \textit{AskReddit}, \textit{politics}, dan \textit{worldnews}. Dataset ini berisi total sekitar 1,3 juta komentar dengan distribusi kelas yang seimbang, yakni 650.000 komentar sarkastis dan 650.000 komentar non-sarkastis \cite{zhang2025enhancing}.

Karakteristik utama dari dataset SARC adalah struktur percakapannya yang bersifat kontekstual. Ekspresi sarkastis dalam Reddit sering kali implisit dan sangat bergantung pada pemahaman terhadap hubungan \textit{parent-child conversation}, berbeda dengan dataset sentimen tradisional yang umumnya hanya mengevaluasi kalimat tunggal secara terisolasi \cite{dubey2025unpacking}. Kompleksitas kontekstual ini, ditambah dengan tingginya penggunaan ungkapan gaul dan struktur tata bahasa yang tidak konsisten, menjadikannya dataset yang ideal untuk mengevaluasi keandalan model pada wacana media sosial yang realistis.

Pemilihan SARC sebagai dataset utama sesuai dengan tujuan utama penelitian ini, terutama dalam konteks analisis trade-off antara kemampuan klasifikasi dan efisiensi komputasi model berbasis transformer. Struktur percakapan yang panjang dan ketergantungan kontekstual dalam diskusi Reddit memberikan tantangan yang cukup kompleks untuk tugas deteksi sarkasme, di mana hubungan semantik jangka panjang perlu ditangkap secara efektif oleh model. Karakteristik ini membuat SARC sangat sesuai untuk menguji kemampuan pemahaman kontekstual RoBERTa, serta untuk mengevaluasi kemampuan DistilBERT dalam mempertahankan kinerja yang kompetitif dengan kebutuhan komputasi yang lebih rendah \cite{dubey2025unpacking}. Selain itu, ukuran dataset yang besar juga memungkinkan analisis yang lebih komprehensif terhadap waktu pelatihan, kecepatan inferensi, penggunaan memori, dan biaya komputasi secara keseluruhan dalam penerapan model.

Meskipun demikian, dataset SARC memiliki sejumlah tantangan terlepas dari keunggulannya. Dataset ini mengandung ekspresi sarkastis yang sangat ambigu yang mungkin bergantung pada makna implisit, referensi budaya, atau konteks percakapan tertentu. Selain itu, sifat komentar yang \textit{\textit{noise}} di Reddit membuat proses preprocessing dan ekstraksi fitur menjadi lebih sulit. Meskipun demikian, tantangan tersebut menjadi bagian penting dalam penelitian ini karena mampu menyediakan benchmark yang realistis untuk mengevaluasi ketangguhan serta perilaku komputasi model transformer dalam mendeteksi sarkasme pada lingkungan media sosial nyata.

\section{Tahap Preprocessing Data}

Tahap preprocessing data dilakukan untuk mempersiapkan dataset bagi model deteksi sarkasme berbasis transformer. Hal ini sangat penting karena teks media sosial dalam \textit{Self-Annotated Reddit Corpus} (SARC) sangat tidak terstruktur dan penuh dengan \textit{noise}. Komentar pada Reddit umumnya mengandung bahasa informal, singkatan, \textit{slang}, kesalahan penulisan, serta berbagai elemen tambahan seperti URL, \textit{user mentions}, dan emoji. Selain itu, ekspresi sarkastis pada Reddit seringkali bersifat implisit dan sangat bergantung pada konteks percakapan sebelumnya, sehingga preprocessing perlu dilakukan secara hati-hati agar informasi kontekstual dan isyarat sarkasme (\textit{sarcasm cues}) tetap dapat dipertahankan.

Pada penelitian ini, preprocessing dilakukan dalam beberapa langkah utama. Tahap awal meliputi pembersihan teks (text cleaning) dengan menghapus elemen non-semantik seperti URL, tag HTML, komentar duplikat, dan user mentions. Sesuai dengan batasan masalah pada pendekatan unimodal teks, elemen multimodal seperti emoji dihapus agar analisis trade-off komputasi pada arsitektur model dapat diukur secara murni tanpa bias komputasi karakter khusus. Selain itu, teks diseragamkan menjadi huruf kecil (lowercasing) dan tanda baca dihapus untuk mengurangi kompleksitas representasi input sesuai standar pra-pemrosesan model bahasa berbasis transformer. Setelah teks dibersihkan, setiap komentar target akan digabungkan dengan komentar induknya (parent comment) untuk menjaga hubungan konteks pada struktur percakapan Reddit. Pendekatan ini penting karena banyak pernyataan sarkastis tidak dapat diinterpretasikan tanpa mempertimbangkan konteks percakapan sebelumnya\cite{dubey2025unpacking, dubey2026efficient}.

Tahap berikutnya adalah melakukan tokenisasi teks dengan tokenizer bawaan dari setiap arsitektur transformer, yaitu RoBERTaTokenizer dan DistilBertTokenizer. Tokenizer bawaan ini digunakan untuk mempertahankan struktur sintaksis dan representasi kontekstual yang diperlukan oleh model transformer untuk menghasilkan embedding kontekstual. Dengan demikian, teknik \textit{pre-processing} tradisional seperti stemming, \textit{aggressive lemmatization}, dan \textit{stopword removal} tidak dilakukan untuk penelitian ini karena berpotensi menghilangkan hubungan semantik antar kata yang sangat penting untuk deteksi sarkasme. Selain itu, penelitian ini juga menggunakan pengaturan maximum \textit{sequence} length tertentu melalui proses truncation untuk mengontrol panjang input teks agar tetap sesuai dengan kapasitas model dan efisien secara komputasi.

Penerapan preprocessing dalam penelitian ini tidak hanya bertujuan untuk meningkatkan kualitas representasi tekstual, tetapi juga untuk mendukung analisis \textit{trade-off} antara performa model dan efisiensi komputasi. Pengurangan \textit{noise} pada teks memungkinkan mekanisme attention untuk lebih fokus pada isyarat sarkasme yang relevan secara linguistik. Di sisi lain, pengurangan panjang sekuens secara optimal dapat membantu mengurangi waktu pelatihan, penggunaan memori GPU, dan kompleksitas komputasi dalam arsitektur transformer \cite{dubey2025unpacking}. Oleh karena itu, fase preprocessing sangat penting dalam memastikan bahwa evaluasi RoBERTa dan DistilBERT dapat dilakukan secara adil, konsisten, dan sejalan dengan tujuan utama penelitian ini.

\section{Implementasi Model RoBERTa}

\begin{figure}[H]
\centering

% Style definitions
\tikzstyle{startstop} = [rectangle, rounded corners,
minimum width=3.2cm,
minimum height=0.8cm,
text centered,
draw=black,
fill=red!20]

\tikzstyle{process} = [rectangle,
minimum width=3.8cm,
minimum height=0.8cm,
text centered,
draw=black,
fill=orange!20]

\tikzstyle{arrow} = [thick,->,>=stealth]

\begin{tikzpicture}[
node distance=1.2cm,
scale=0.78,
transform shape
]

% Nodes
\node (input) [startstop]
{Parent Comment + Target Comment};

\node (preprocess) [process, below of=input]
{Text Preprocessing};

\node (tokenizer) [process, below of=preprocess]
{RoBERTa Tokenizer};

\node (embedding) [process, below of=tokenizer]
{Input Embedding Layer};

\node (encoder) [process, below of=embedding]
{\shortstack{Transformer Encoder \\ (Multi-Head Self-Attention)}};

\node (contextual) [process, below of=encoder]
{Contextual Representation};

\node (classification) [process, below of=contextual]
{Classification Layer};

\node (output) [startstop, below of=classification]
{Sarcastic / Non-Sarcastic Output};

% Arrows
\draw [arrow] (input) -- (preprocess);
\draw [arrow] (preprocess) -- (tokenizer);
\draw [arrow] (tokenizer) -- (embedding);
\draw [arrow] (embedding) -- (encoder);
\draw [arrow] (encoder) -- (contextual);
\draw [arrow] (contextual) -- (classification);
\draw [arrow] (classification) -- (output);

\end{tikzpicture}

\caption{Architecture of RoBERTa implementation for sarcasm detection on the SARC dataset}
\label{fig:robertaflow}

\end{figure}

Pada penelitian ini, RoBERTa digunakan sebagai salah satu model berbasis transformer utama untuk deteksi sarkasme pada dataset \textit{Self-Annotated Reddit Corpus} (SARC). RoBERTa \textit{(Robustly Optimized BERT Pretraining Approach)} adalah versi modifikasi dari arsitektur BERT yang dioptimalkan untuk pemahaman konteks bahasa yang lebih baik melalui \textit{pre-training} yang lebih optimal. RoBERTa memperkenalkan beberapa peningkatan arsitektural dibandingkan model BERT standar, termasuk \textit{dynamic masking}, penghapusan objektif \textit{Next Sentence Prediction} (NSP), dan pelatihan dataset yang lebih besar untuk jangka waktu yang lebih lama. Peningkatan tersebut memungkinkan RoBERTa menghasilkan representasi kontekstual yang lebih kaya dan mampu menangkap hubungan semantik antar kata dengan lebih baik.

Penelitian ini menggunakan model RoBERTa karena sifat sarkasme yang bergantung pada konteks dalam percakapan media sosial. Ucapan sarkastis biasanya dipahami sebagai kebalikan dari apa yang diucapkan secara harfiah dan seringkali memerlukan konteks percakapan sebelumnya agar dapat diinterpretasikan dengan benar. Model RoBERTa menggunakan mekanisme \textit{multi-head self-attention} yang memungkinkan model untuk menganalisis hubungan antar kata dan petunjuk konteks secara bersamaan, sehingga sangat sesuai untuk mengidentifikasi sarkasme implisit pada diskusi Reddit. Selain itu, kemampuan \textit{contextual embedding} RoBERTa memungkinkan model untuk menangani teks media sosial yang \textit{noise} dan informal dengan lebih baik, seperti bahasa gaul, singkatan, dan struktur tata bahasa yang tidak konsisten, yang sering ditemukan dalam dataset SARC \cite{dubey2025unpacking, jimenez2026bert}.

Proses implementasi dimulai dengan menggunakan RoBERTaTokenizer untuk melakukan tokenisasi teks hasil preprocessing sebelumnya. Setiap input terdiri dari komentar target yang digabungkan dengan \textit{parent comment} untuk mempertahankan konteks percakapan dalam struktur diskusi Reddit. Hasil tokenisasi kemudian dikonversi menjadi \textit{contextual embedding} melalui beberapa lapisan encoder transformer. Pada tahap fine-tuning, model RoBERTa yang telah dilatih sebelumnya disesuaikan dengan tugas deteksi sarkasme dengan data berlabel dari dataset SARC. Representasi kontekstual terakhir dari encoder transformer dilewatkan melalui lapisan klasifikasi dengan fungsi aktivasi Softmax untuk memprediksi kelas komentar (sarkastis atau non-sarkastis).

Meskipun RoBERTa memiliki pemahaman kontekstual yang sangat baik, arsitektur ini menghadirkan tantangan biaya komputasi yang substansial. Secara bawaan, mekanisme \textit{self-attention} pada \textit{transformer} memiliki tingkat kompleksitas yang tinggi, yang menjadi semakin membebani ketika memproses sekuens teks panjang dari kumpulan data berskala besar seperti SARC. Akibatnya, RoBERTa menuntut kapasitas memori GPU yang masif, waktu pelatihan yang lebih lambat, serta menghasilkan latensi inferensi yang jauh lebih tinggi dibandingkan model \textit{transformer} ringan (\textit{lightweight}) \cite{dubey2025unpacking}. Namun demikian, karakteristik komputasi yang berat inilah yang justru menjadikan RoBERTa sangat relevan dengan tujuan penelitian ini, yaitu sebagai representasi model berkinerja tinggi untuk menganalisis \textit{trade-off} antara akurasi klasifikasi dan efisiensi komputasi pada sistem deteksi sarkasme.

\section{Implementasi Model DistilBERT}

\begin{figure}[H]
\centering

% Style definitions
\tikzstyle{startstop} = [rectangle, rounded corners,
minimum width=3.2cm,
minimum height=0.8cm,
text centered,
draw=black,
fill=red!20]

\tikzstyle{process} = [rectangle,
minimum width=3.8cm,
minimum height=0.8cm,
text centered,
draw=black,
fill=orange!20]

\tikzstyle{arrow} = [thick,->,>=stealth]

\begin{tikzpicture}[
node distance=1.2cm,
scale=0.78,
transform shape
]

% Nodes
\node (input) [startstop]
{Parent Comment + Target Comment};

\node (preprocess) [process, below of=input]
{Text Preprocessing};

\node (tokenizer) [process, below of=preprocess]
{DistilBERT Tokenizer};

\node (distillation) [process, below of=tokenizer]
{Knowledge Distillation};

\node (embedding) [process, below of=distillation]
{Compressed Embedding Layer};

\node (encoder) [process, below of=embedding]
{Compressed Transformer Encoder};

\node (contextual) [process, below of=encoder]
{Contextual Representation};

\node (classification) [process, below of=contextual]
{Classification Layer};

\node (output) [startstop, below of=classification]
{Sarcastic / Non-Sarcastic Output};

% Arrows
\draw [arrow] (input) -- (preprocess);
\draw [arrow] (preprocess) -- (tokenizer);
\draw [arrow] (tokenizer) -- (distillation);
\draw [arrow] (distillation) -- (embedding);
\draw [arrow] (embedding) -- (encoder);
\draw [arrow] (encoder) -- (contextual);
\draw [arrow] (contextual) -- (classification);
\draw [arrow] (classification) -- (output);

\end{tikzpicture}

\caption{Architecture of DistilBERT implementation for sarcasm detection on the SARC dataset}
\label{fig:distilbertflow}

\end{figure}

Pada penelitian ini, kami menggunakan DistilBERT sebagai model pembanding berbasis transformer untuk menganalisis trade-off antara performa klasifikasi dan efisiensi komputasi dalam tugas pendeteksian sarkasme. DistilBERT merupakan versi ringan (lightweight transformer) dari arsitektur BERT yang dikembangkan menggunakan teknik knowledge distillation. Pendekatan ini memungkinkan model untuk mewarisi pengetahuan dari arsitektur transformer yang lebih besar, sambil mempertahankan sebagian besar kemampuannya untuk memahami konteks. DistilBERT mengurangi sekitar 40\% lapisan yang digunakan dalam arsitektur BERT standar. Hal ini menghasilkan ukuran parameter yang lebih kecil dan kompleksitas komputasi yang lebih rendah.

Penggunaan DistilBERT didasarkan pada kebutuhan sistem deteksi sarkasme yang tidak hanya kompetitif secara performa, tetapi juga adaptif secara operasional. Pada platform media sosial seperti Reddit, sistem berpotensi diterapkan pada skenario \textit{real-time} dengan volume data masif yang menuntut latensi inferensi sangat rendah. Sebagai alternatif dari model berat seperti RoBERTa, DistilBERT menawarkan kecepatan pelatihan dan inferensi yang jauh lebih tinggi. Meskipun mengalami kompresi arsitektural yang signifikan, DistilBERT terbukti secara empiris mampu mempertahankan kapasitas pemahaman kontekstual yang memadai untuk menangkap ekspresi sarkasme yang implisit di dalam wacana sosial \cite{dubey2025unpacking}.

Proses implementasi dimulai dengan tahap tokenisasi teks yang sudah diproses sebelumnya menggunakan DistilBertTokenizer. Setiap input terdiri dari komentar target yang digabungkan dengan komentar induknya (parent comment) untuk mempertahankan konteks percakapan Reddit, seperti pada implementasi RoBERTa. Hasil tokenisasi kemudian diubah representasi contextual embeddings melalui beberapa lapisan transformer encoder yang terkompresi. Pada tahap fine-tuning, model DistilBERT disesuaikan dengan tugas klasifikasi sarkasme menggunakan dataset SARC yang telah diberi label. Representasi kontekstual terakhir kemudian diteruskan melalui lapisan klasifikasi dengan fungsi aktivasi Softmax untuk memprediksi kelas sarkastis dan non-sarkastis.

Keuntungan utama dari DistilBERT adalah efisiensi komputasinya. Parameter yang lebih sedikit memungkinkan model ini menggunakan memori GPU yang jauh lebih rendah, mempercepat waktu pelatihan, dan memangkas latensi inferensi dibandingkan dengan RoBERTa. Karakteristik ini membuat DistilBERT sangat ideal untuk diterapkan di lingkungan dengan keterbatasan sumber daya maupun kebutuhan pemrosesan \textit{real-time}. Meskipun demikian, arsitektur DistilBERT yang dikompresi juga berdampak pada berkurangnya kedalaman representasi kontekstual. Hal ini menyebabkan performa klasifikasinya sedikit lebih rendah dibandingkan RoBERTa, terutama ketika berhadapan dengan kasus sarkasme yang sangat implisit dan kompleks \cite{dubey2025unpacking}. Oleh karena itu, penerapan DistilBERT dalam penelitian ini menjadi krusial untuk mengevaluasi sejauh mana efisiensi komputasi dapat dicapai tanpa menimbulkan degradasi yang signifikan pada performa deteksi sarkasme..

\section{Metode Evaluasi dan Pengukuran Kinerja}

Dalam penelitian ini, model RoBERTa dan DistilBERT untuk deteksi sarkasme dievaluasi dalam hal kinerja klasifikasi dan efisiensi komputasi menggunakan dataset SARC. Proses evaluasi tidak hanya mempertimbangkan kemampuan model dalam mengklasifikasikan komentar sarkastis dan non-sarkastis secara akurat, tetapi juga sumber daya komputasi yang diperlukan untuk pelatihan dan inferensi. Sebagai hasilnya, penelitian ini menggunakan dua kategori utama evaluasi, yaitu metrik kinerja klasifikasi dan metrik efisiensi komputasi.

\subsection{Classification Performance Metrics}

Evaluasi performa klasifikasi dilakukan untuk mengetahui kemampuan model RoBERTa dan DistilBERT dalam membedakan komentar sarkastis dan non-sarkastis pada dataset pengujian. Untuk memperoleh gambaran performa yang lebih menyeluruh, penelitian ini menggunakan metrik \textit{accuracy}, \textit{precision}, \textit{recall}, dan \textit{F1-score}.

Masing-masing metrik memberikan sudut pandang yang berbeda dalam proses evaluasi. \textit{Accuracy} digunakan untuk melihat tingkat ketepatan prediksi model secara keseluruhan, sedangkan \textit{precision} dan \textit{recall} digunakan untuk menilai kemampuan model dalam mengidentifikasi komentar yang mengandung sarkasme. Selain itu, \textit{F1-score} digunakan sebagai metrik utama karena mampu merepresentasikan keseimbangan antara \textit{precision} dan \textit{recall}, sehingga dapat memberikan gambaran yang lebih objektif terhadap performa model.

Perhitungan dan definisi masing-masing metrik mengacu pada landasan teori yang telah dijelaskan pada Subbab 2.3.

\subsection{Computational Efficiency Metrics}

Selain mengevaluasi performa klasifikasi, penelitian ini juga menilai efisiensi komputasi dari setiap model yang digunakan. Evaluasi tersebut dilakukan untuk mengetahui besarnya sumber daya komputasi yang dibutuhkan selama proses pelatihan maupun inferensi, sehingga dapat memberikan gambaran mengenai kelayakan penerapan model pada kondisi nyata.

Pengukuran efisiensi komputasi dilakukan menggunakan beberapa metrik, yaitu \textit{training time}, \textit{inference latency}, penggunaan memori GPU, dan jumlah parameter model. \textit{Training time} digunakan untuk mengukur waktu yang diperlukan selama proses \textit{fine-tuning}, sedangkan \textit{inference latency} digunakan untuk melihat kecepatan model dalam menghasilkan prediksi pada data pengujian. Penggunaan memori GPU diamati untuk mengetahui kebutuhan sumber daya selama proses pelatihan, sementara jumlah parameter model digunakan sebagai indikator kompleksitas arsitektur yang dimiliki oleh masing-masing model.

Definisi dan dasar teori dari metrik efisiensi komputasi yang digunakan dalam penelitian ini mengacu pada Subbab 2.3.

\subsection{Trade-off Analysis}

Tahap evaluasi akhir dilakukan melalui analisis trade-off antara performa klasifikasi dan efisiensi komputasi dari RoBERTa dan DistilBERT. Analisis ini tidak bertujuan untuk menentukan satu model yang terbaik secara mutlak, tetapi untuk melihat seberapa besar perubahan performa yang terjadi dibandingkan dengan keuntungan efisiensi komputasi yang diperoleh.

Performa klasifikasi dibandingkan berdasarkan accuracy, precision, recall, dan F1-score, sedangkan efisiensi komputasi dibandingkan berdasarkan training time, inference latency, penggunaan memori GPU, dan jumlah parameter model. Perbedaan tersebut kemudian dianalisis dalam bentuk penurunan performa (performance loss) dan peningkatan efisiensi (efficiency gain), termasuk speedup factor untuk melihat perbedaan kecepatan antara kedua model.

Hasil analisis digunakan untuk melihat kondisi penggunaan masing-masing model. RoBERTa dapat lebih sesuai ketika performa klasifikasi menjadi prioritas, sedangkan DistilBERT dapat lebih sesuai ketika efisiensi komputasi dan keterbatasan sumber daya menjadi pertimbangan utama.


\section{Skenario Pengujian dan Analisis Trade-off}

Dalam penelitian ini, skenario pengujian dirancang untuk membandingkan performa klasifikasi dan efisiensi komputasi RoBERTa dan DistilBERT pada dataset SARC. Agar perbandingan dilakukan secara adil, kedua model menggunakan dataset, preprocessing, struktur input parent comment dan target comment, serta lingkungan komputasi yang sama.

Dataset yang digunakan berupa sekitar 200.000 pasangan parent comment–target comment yang dipilih secara terkontrol dari SARC. Data kemudian dibagi secara stratified menjadi 80\% data pelatihan, 10\% data validasi, dan 10\% data pengujian. Panjang maksimum input ditetapkan sebesar 256 token. Sebelum eksperimen utama dilakukan, pengujian awal menggunakan sekitar 20.000–50.000 data dilakukan untuk memastikan pipeline berjalan dengan baik serta menentukan batch size yang sesuai dengan kapasitas GPU.

Eksperimen utama dilakukan menggunakan GPU NVIDIA RTX 5070 Ti dengan VRAM 16 GB yang tersedia di AI Lab. Proses pelatihan menggunakan mixed precision atau FP16 untuk menyesuaikan kebutuhan komputasi dengan kapasitas perangkat keras. Google Colab dan perangkat pribadi hanya digunakan untuk tahap pengembangan, preprocessing, serta pengujian awal, sedangkan seluruh pengukuran akhir dilakukan pada perangkat GPU yang sama agar hasil perbandingan tetap konsisten.

Setiap model dilatih sebanyak tiga kali menggunakan random seed yang sama untuk mengurangi kemungkinan hasil dipengaruhi oleh satu kali proses pelatihan. Konfigurasi fine-tuning seperti learning rate, batch size, dan jumlah epoch dapat disesuaikan dengan karakteristik masing-masing model, tetapi proses penentuannya dilakukan dengan prosedur yang sama agar perbandingan tetap adil.

Performa klasifikasi dievaluasi menggunakan accuracy, precision, recall, dan F1-score. Efisiensi komputasi diukur berdasarkan training time, inference latency, penggunaan memori GPU, dan jumlah parameter model. Hasil dari ketiga pengujian kemudian dirangkum untuk melihat konsistensi performa dan kebutuhan komputasi masing-masing model.

Analisis trade-off dilakukan dengan membandingkan perubahan performa klasifikasi terhadap keuntungan efisiensi komputasi yang diperoleh. Selain itu, perbedaan kecepatan antara kedua model dihitung menggunakan speedup factor \cite{dubey2025unpacking}:

\[
\text{Speedup Factor} =
\frac{T_{\text{RoBERTa}}}{T_{\text{DistilBERT}}}
\]

Pengukuran tersebut dilakukan secara terpisah untuk waktu pelatihan dan latensi inferensi. Hasil akhir digunakan untuk melihat kondisi penggunaan masing-masing model, yaitu ketika performa klasifikasi menjadi prioritas utama, ketika keterbatasan sumber daya menjadi pertimbangan utama, maupun ketika diperlukan keseimbangan antara keduanya.


%

%%%%%%%%%%%%%%%%%%%%%%%%%%%%%%%%%%%%%%%%%%%%%%%%%%%%%%%%%%%%%%%
% Berikut adalah untuk membuat Daftar pustaka
% Style daftar pustaka menggunakan format IEEE
%%%%%%%%%%%%%%%%%%%%%%%%%%%%%%%%%%%%%%%%%%%%%%%%%%%%%%%%%%%%%%%

\cleardoublepage
\phantomsection
\addcontentsline{toc}{chapter}{Daftar Pustaka}

\bibliographystyle{IEEEtran}
\bibliography{References}

%%%%%%%%%%%%%%%%%%%%%%%%%%%%%%%%%%%%%%%%%%%%%%%%%%%%%%%%%%%%%%%
% Berikut untuk lampiran
%%%%%%%%%%%%%%%%%%%%%%%%%%%%%%%%%%%%%%%%%%%%%%%%%%%%%%%%%%%%%%%

\cleardoublepage
\addcontentsline{toc}{chapter}{Lampiran}
\include{Lampiran}

\end{document}
